% RESEARCH_DRAFT_SHA256: e7e170784bd760a27a174629782b1cbfb91f21b84e158f213debf224b406efae
% Layout and source-reference wrappers expanded for the web renderer.
\documentclass[11pt]{article}
\usepackage[margin=25mm]{geometry}
\usepackage{fontspec,unicode-math}
\setmainfont{Latin Modern Roman}
\setmathfont{Latin Modern Math}
\usepackage{amsmath,amssymb,mathtools}
\usepackage{xurl,hyperref}
\hypersetup{colorlinks=true,urlcolor=blue,linkcolor=blue}
\setcounter{secnumdepth}{0}
\setlength{\parindent}{0pt}
\setlength{\parskip}{5pt}
\setlength{\emergencystretch}{3em}
\newcommand{\sref}[2]{\hyperlink{source:#1:#2}{[S#2]}}
\newcommand{\cataloguescope}[1]{\paragraph{Source question.}#1}
\newcommand{\floor}[1]{\left\lfloor#1\right\rfloor}
\begin{document}
\section*{Research-attempt provenance}
This standalone research attempt retains the complete original definition body below,
including its historical source and preparation notes. Those notes describe the
earlier catalogue preparation; the new attempt and its verification are identified
separately after that body. No original ID, catalogue status, or source statement is changed.

% UnsolvedMath ID 30003993; source catalogue code OWR-16633-010
\section{Integrality Gaps for Directed Linear k-Cut}\label{30003993}
{\small UnsolvedMath ID 30003993 · Computer Science}
\subsection{Definitions and mathematical statement}

\paragraph{Shared notation.}Standard number systems are denoted by $\mathbb N,\mathbb Z,\mathbb Q,\mathbb R,\mathbb C$; index ranges, computational models and specialized notation follow the cited problem source and the definitions below.


Directed Linear-$k$-Cut specifies ordered terminals and removes arcs to block paths in the prescribed forward-terminal directions. Its natural path relaxation gives fractional arc deletion variables satisfying the terminal-separation constraints.

A relaxation contains the integral feasible solutions. For a nonnegative minimization objective, its integrality gap is the supremum of the ratio between the integral optimum and the relaxation optimum over the specified instances.

A linear program optimizes $c^Tx$ over a polyhedron $\{x:Ax\leq b\}$. In the simplex method a pivot moves between adjacent basic feasible solutions; a pivot rule chooses the next such move.
\paragraph{Statement recorded in the supplied source.}
For directed Linear-$k$-Cut, is the integrality gap of the natural linear-programming relaxation $O(1)$ for every $k$, or is it $\omega(1)$? \href{https://doi.org/10.4171/owr/2018/50}{[S2]}
\paragraph{Scope qualification.}
The exact ordered-pair convention and LP formulation are those of the report.
\paragraph{Source status.}
The supplied catalogue labels this record open. Its source-status metadata was checked against the supplied snapshot on 2026-10-11. This import does not independently establish current open status or mathematical correctness. The catalogue literature assessment is dated 2026-08-21.
\subsection{Short English statement}
Decide whether the directed ordered-cut relaxation has a uniform constant gap or an unbounded gap.
\subsection{Sources}
\begin{itemize}
\item[S1] Ulam AI and the UnsolvedMath Contributors, UnsolvedMath: A Curated Collection of Open Mathematics Problems, record 30003993 (OWR-16633-010). Catalogue curation and metadata: CC BY 4.0. The supplied source question is retained; definitions and the English summary are editorial additions.
\url{https://huggingface.co/datasets/ulamai/UnsolvedMath}.
\item[S2] Combinatorial Optimization (2019), Oberwolfach Reports, DOI 10.4171/owr/2018/50, pp. 2969--3023
\url{https://doi.org/10.4171/owr/2018/50}.
\item[S3] Linear 3-Cut and General Linear k-Cut
\url{https://karthik.ise.illinois.edu/pubs/linear-3-cut.pdf}.
\textit{Underlying Oberwolfach Reports material retains its source-specific licence, including CC BY-SA 4.0 where applicable; see THIRD\_PARTY\_NOTICES.md.}
\end{itemize}
\paragraph{Preparation and review.}
Original question attributed to the cited source and UnsolvedMath contributors. Codex (GPT-6; exact serving revision not exposed) assisted with definitions, English summaries, source inspection and formatting. No new human mathematical authorship or independent proof verification is asserted. Canonical ID 30003993 and source code OWR-16633-010 checked against the complete supplied registry and existing repository statements; known flagged repetitions and obvious identical conjectures excluded. Full mathematical-equivalence/literature review has not been performed. Statement preparation and final review followed contribution prompts C and D; automated checks verify structure and identity, not open status.
\label{end:30003993}

\section{Literature and context check}
The report, p. 3013, asks whether the general Linear-k-Cut gap is bounded independently of k. The consulted Berczi--Chandrasekaran--Kiraly--Madan paper distinguishes edge and node variants and records their three-terminal results. This attempt works directly with the edge-path LP and reproves a logarithmic route already noted by the report; it does not claim a new approximation guarantee.
\section{Research attempt: Recursive threshold rounding and the missing constant bound}
Let edge costs $c_e\ge0$, terminals $s_1,\ldots,s_k$, and an LP solution $x_e\ge0$ satisfy $\sum_{e\in P}x_e\ge1$ for every directed $s_i$-to-$s_j$ path with $i<j$. The LP objective is $L=\sum_ec_ex_e$. Unit costs give the report's edge-count setting.

Split the ordered terminal list into a left and right half. Let $d(v)$ be the shortest $x$-length from any left terminal to $v$, capped at one, taking value one if unreachable. Left terminals have distance zero, and right terminals have distance one by the path inequalities. For an arc $u\to v$, $d(v)\le d(u)+x_{uv}$; this follows from extending a shortest path when $d(u)<1$, and is automatic when $d(u)=1$. Choose $\theta$ uniformly in $(0,1)$, put $U=\{v:d(v)<\theta\}$, and delete all arcs from $U$ to its complement. Such an arc is deleted with probability at most $\max(0,d(v)-d(u))\le x_{uv}$. The expected deleted cost is at most $L$ by summing these probabilities.

Now recurse separately on the two induced vertex sets, using the inherited ordered terminal lists and inherited LP values. Every restricted path was an original path, so the LP remains feasible inside each subproblem. At one recursion level, the vertex sets are pairwise disjoint; consequently the sums of their inherited $c_ex_e$ are at most the original $L$. Conditional on every earlier partition, the expected new deletion cost on that level is at most $L$. The recursion has at most $\lceil\log_2 k\rceil$ nontrivial levels, giving expected total cost at most $\lceil\log_2 k\rceil L$ for $k\ge2$.

The final deletion is feasible even though backward arcs between regions remain. Each final region contains one terminal, and regions inherit their terminal order. For any remaining arc joining two different final regions, examine the first split that separated its endpoints. Arcs from the left part to the right part were deleted at that split, so this remaining arc goes from the higher-ordered region to the lower-ordered one. Thus region indices are nonincreasing along every remaining path, making a path from $s_i$ to $s_j$ with $i<j$ impossible. For $k\le1$, the empty cut suffices.

This proves the known logarithmic upper bound by a complete route. Removing its logarithm needs a new argument: one conditional estimate per level cannot alone imply a constant total. As a precise abstract obstruction, $h$ independent events of probability $x=1/h$ have union probability $1-(1-1/h)^h$. Since $1-t\le e^{-t}$, this is at least $1-e^{-1}$, whose ratio to $x$ grows linearly with $h$. Thus per-level marginal bounds of size $x$ permit logarithmic accumulation. These events are not claimed to be an actual Linear-k-Cut integrality-gap instance; they show exactly why that information alone cannot certify a constant bound.
\section{Verification and adversarial audit}
The rounding proof checks terminal inclusion, capped infinite distances, induced-subproblem LP feasibility, and paths that might leave a recursive region. Region-order monotonicity rules out such circumvention. The event example is explicitly a limitation of the analysis, not a construction proving an unbounded LP gap.
\paragraph{Audit corrections.} The LP is explicitly defined, and the report's unit-cost case is included. A logarithmic rounding analysis is not advertised as a lower bound on the true gap.
\paragraph{External computation.} None; symbolic arguments and boundary checks only.
\section{Scope and remaining gaps}
The original constant-versus-unbounded gap dichotomy remains unresolved. The proof rederives a known O(log k) bound and locates the failure of a naive constant-factor amortization.
\paragraph{Integration review.} Codex (GPT-6; exact serving revision not exposed) reviewed the written arguments and expanded layout and source-reference wrappers. Repository metadata and the production web renderer were checked using Python 3.13 and JavaScript with MathJax 3.2.2. This model-assisted review is not independent human proof verification; no PDF compilation or formal proof checking is asserted.

\section{Completion Estimate}
\noindent\textbf{COMPLETION: 0\%}
The calculations below the preserved source statement investigate a concrete route,
but establish no substantiated new progress toward the entire original target.
Known special cases and failed sufficient conditions are not credited as a solution.
\section{References}
The original source references and catalogue attribution are preserved in the
Sources section above. The following URLs were actually consulted for this attempt;
retention of an additional original reference does not claim fresh consultation of it.
\begin{itemize}
\item \url{https://ems.press/content/serial-article-files/46772} (consulted 2026-10-11). 
\item \url{https://karthik.ise.illinois.edu/pubs/linear-3-cut.pdf} (consulted 2026-10-11). 
\end{itemize}
\section{Model and contributor disclosure}
OpenAI gpt-6-astra, requested reasoning configuration ultra, generated this research
attempt on 2026-10-11. The exact serving revision is not exposed. Codex GPT-6
performed orchestration, repository-format review, and adversarial mathematical review of the written argument; its exact serving revision
is also not exposed. The mathematical route, boundary checks and final adversarial
self-audit followed contribution prompts A then B in
\texttt{docs/contribution\_prompts.md}, after reading \texttt{CONTRIBUTING.md}.
Python 3.14.2 assembled the standalone documents; any mathematical computation is
specified in Verification. No theorem prover was used. No human mathematical
authorship, substantial human proof correction, or independent human proof checking
is asserted. Model-assisted review, source retrieval, and structural checks do not
establish correctness or novelty. Full private reasoning transcripts are unavailable
as reproducibility artifacts; the complete arguments, inputs and outcomes actually
used are stated here or in the companion batch audit.
\section{Final status}
Unresolved as a new research contribution. This is a concrete route investigation,
not a claim of new complete resolution or a current-status certification. A known
answer reproduced for context does not become a new contribution. The catalogue's
source-status label is preserved; this attempt status does not assert that a
source-labelled solved question is open.
\end{document}
